\subsection{The marked handle evaluator}
\label{subsec:graphical-construction}

\begin{definition}[Decorated Kirby ribbon graph]
\label{def:decorated-kirby-graph}
\label{def:dotted-circle-translation}
\label{def:colored-graph-sum}
Fix a twisted categorical input \((\cC,\eta)\).  Let \(\mathcal H\) be a
one-ended Kirby presentation, let \(\omega\in Z^2(\mathcal H)\), and assign a
simple object \(\lambda_k\in\operatorname{Irr}(\cC)\) to every framed
component \(K_k\).  Choose auxiliary orientations of the framed components,
write \(e_k\) for the corresponding oriented \(2\)-handle generator, and put
\[
  \omega_k:=\omega(e_k)\in\mathbb Z.
\]
Choose also an admissible section system for the dotted components.  We write
\[
  L^\eta_{\lambda,\omega}(\mathcal H)
\]
for the following closed coloured ribbon graph, understood as a finite linear
combination whenever a projector is expanded.  The component \(K_k\) is
coloured by \(\lambda_k\) and carries the coupon
\[
  (\eta^{\omega_k})_{\lambda_k}
  =\eta_{\lambda_k}^{\omega_k}\operatorname{id}_{\lambda_k}.
\]
Here \(\eta\colon\operatorname{Id}_{\cC}\Rightarrow
\operatorname{Id}_{\cC}\) is the character in
\cref{def:twisted-categorical-input}; \(\eta^{\omega_k}\) is its integer
power, using the inverse when \(\omega_k<0\); and
\(\eta_{\lambda_k}\in\mathbb C^\times\) is the simple-object scalar of
\cref{lem:simple-object-character-scalars}.  Reversing the auxiliary
orientation of \(K_k\) replaces \(e_k,\omega_k,\lambda_k\) by
\(-e_k,-\omega_k,\lambda_k^*\); this does not change the evaluator by
\cref{prop:orientation-choice-independence}.

A dotted component \(\delta\) is not coloured.  Instead, cut its handle as in
\cref{def:cut-open-dotted-tangle}.  If \(B_\delta\) is the tensor product
specified by the ordered oriented strands on one cut face, insert the
canonical projector
\[
  P_{\mathbf1,B_\delta}\colon B_\delta\longrightarrow B_\delta
\]
and reglue the two faces.  Doing this at every dotted component gives
\(L^\eta_{\lambda,\omega}(\mathcal H)\).  The Kirby-colour coefficient
\(\prod_k d_{\lambda_k}\) is not included in this notation.
\end{definition}

\begin{definition}[Homologically twisted Crane--Yetter evaluator]
\label{def:twisted-handle-evaluator}
Fix a twisted categorical input \((\cC,\eta)\), a one-ended Kirby presentation
\(\mathcal H\), and a cocycle \(\omega\in Z^2(\mathcal H)\).  Then the handle
evaluator is
\begin{equation}
  \boxed{
  \CY_{\cC,\eta}(\mathcal H,\omega)
  :=
  \operatorname{Dim}(\cC)^{1-h_3(\mathcal H)}
  \sum_{\lambda}
    \left(\prod_k d_{\lambda_k}\right)
    \operatorname{RT}_{\cC}
      \bigl(S^3,L^\eta_{\lambda,\omega}(\mathcal H)\bigr).}
  \label{eq:twisted-handle-evaluator}
\end{equation}
Here the sum runs over all assignments of simple objects to the framed
components; \(d_{\lambda_k}\) is the quantum dimension of the assigned
simple object; \(\operatorname{Dim}(\cC)=\sum_i d_i^2\) is the global
dimension; \(h_3(\mathcal H)\) is the number of \(3\)-handles; and
\(\operatorname{RT}_{\cC}(S^3,-)\) is the closed ribbon-graph evaluation of
\cref{def:rt-graph-evaluation}.
\end{definition}

\begin{remark}[Three-dimensional meaning]
\label{rem:planar-projection-after-cutting}
The dotted--framed link and the decorated graph of
\cref{def:decorated-kirby-graph} live in the standard
\(S^3=\partial B^4\).  To draw them, choose a point at infinity disjoint from
the graph and identify its complement with \(\mathbb R^3\), then choose a
generic projection \(\mathbb R^3\to\mathbb R^2\).  Changing either choice
changes only the ribbon diagram, not the embedded coloured ribbon graph.
Ribbon graphical coherence therefore gives the same RT evaluation.
\end{remark}

\begin{proposition}[Dependence on the three-handle data]
\label{prop:evaluator-three-handle-dependence}
Although a one-ended Kirby presentation contains the full attaching embeddings
of its \(3\)-handles, the marked evaluator uses that data only through
\[
  h_3(\mathcal H)
  \quad\text{and}\quad
  \operatorname{im}(\partial_3)\subset C_2(\mathcal H).
\]
More precisely, the scalar in
\eqref{eq:twisted-handle-evaluator} reads \(h_3(\mathcal H)\) in its
normalization, while the requirement that
\(\omega\in Z^2(\mathcal H)\) depends on \(\partial_3\) only through the
condition that \(\omega\) vanish on \(\operatorname{im}(\partial_3)\).
\end{proposition}

\begin{proof}
The coloured ribbon graph and its label sum are constructed entirely from
the dotted--framed \(1\)- and \(2\)-handle diagram.  The displayed formula
contains the \(3\)-handles explicitly only in the exponent
\(1-h_3(\mathcal H)\).  Finally,
\[
  \omega\in Z^2(\mathcal H)
  \quad\Longleftrightarrow\quad
  \omega\circ\partial_3=0,
\]
which depends only on the image of \(\partial_3\).  No other part of the
attaching embeddings enters the formula.
\end{proof}

\begin{proposition}[Finiteness of the evaluator]
\label{prop:evaluator-finiteness}
For a one-ended Kirby datum, the sum in
\eqref{eq:twisted-handle-evaluator} is finite.
\end{proposition}

\begin{proof}
The Kirby-colour sums are finite because \(\cC\) has finitely many simple
isomorphism classes.  Each tensor-unit projector is a finite sum because
\(\cC\) is finite semisimple.  A Kirby diagram has finitely many components,
so only finitely many choices occur.
\end{proof}

\begin{remark}[What the dotted circle is not]
\label{rem:dotted-circle-operation-scope}
A dotted circle is not a coloured ribbon component and is never assigned the
Kirby colour.  It denotes the projector onto the tensor-unit through-channel.
In a modular category that projector admits an alternative encircling
description using \(\operatorname{Dim}(\cC)^{-1}\Omega_{\cC}\), but no such
replacement is used in the premodular construction.
\end{remark}

\begin{remark}[Intrinsic meaning]
\label{rem:graph-sum-intrinsic-meaning}
The expression \(L^\eta_{\lambda,\omega}(\mathcal H)\) denotes the resulting
intrinsic composite modulo ribbon graphical relations, not a preferred planar
drawing.  In particular, the orientation of a ribbon determines whether its
label is \(X\) or \(X^*\); its passage through a dotted-handle cross-section
does not recolour it.
\end{remark}
